% Unnumbered research record. No new A5 result or qualifying movement.
\section*{The exact missing finite-scaffold bound}

This records the remaining step after the finite-pair contraction test.
It does not assert that a constrained homeomorphism pool exists.

\textbf{What finite pair types do supply.}
For two fixed definable families of compact sets with finitely many
closed marked subsets, the relation that two members are definably
homeomorphic preserving the marks has a finite formula certificate pool.
The assertion holds in every model of the complete theory with the
coefficients of the two families named.

Here is a proof that also specifies the all-model uniformity. In the
original model, simultaneous Hardt triviality and a triangulation of one
marked reference fibre for each piece give finitely many standard
rational finite polyhedral systems $(K_i;K_i^1,\ldots,K_i^r)$ and graph
formulas for marked homeomorphisms
\[
 \alpha_{i,a}:K_i\longrightarrow P_a,\qquad
 \beta_{j,b}:L_j\longrightarrow Q_b.
 \tag{SF.1}
\]
Their partition pieces and homeomorphisms have finitely many additional
coefficients $d$. The validity of this finite system---the partitions
cover their parameter spaces, each indicated map and inverse is
continuous, and all the marked-set equalities hold---is one first-order
condition $V(d)$. The sentence $\exists d\,V(d)$ holds in the original
model, hence in every model of the fixed complete theory. This argument
does not assume that one particular coefficient tuple belongs to all
models.

For each of the finitely many pairs $(K_i,L_j)$ label whether there is a
marked PL homeomorphism. This label is invariant over real closed
fields. Indeed, a definable marked homeomorphism gives a marked PL one
by relative definable Hauptvermutung; its actual finite subdivisions and
affine pieces give an ordered-field certificate, so existence reflects
to the real algebraic field. Conversely, a chosen real-algebraic PL
homeomorphism is a definable witness in every model. The precise original
input is Shiota, arXiv:1002.1508v3, Theorem~2.1, printed p.~5, as already
audited for A5.019.

Choose one such map $h_{ij}$ for each positive label. Every required
marked homeomorphism is then witnessed by one of the finitely many
formulas
\[
              \beta_{j,b}\,h_{ij}\,\alpha_{i,a}^{-1}.
 \tag{SF.2}
\]
The extra tuple $d$ is part of the certificate. If a marked homeomorphism
between $P_a,Q_b$ exists, transport it through any valid system (SF.1);
the corresponding pair must have a positive label, so (SF.2) supplies a
witness. The reverse implication is immediate. Joint continuity in
$(a,b)$ across the Hardt pieces is not claimed.

\textbf{The constraint which is not supplied by that proof.}
Already for fixed compact polyhedra $P,Q$ and a fixed definable family
of embeddings $r_w:P\to R^N$, $e_w:Q\to R^N$, the controlled question is
\[
 \exists h:P\xrightarrow{\cong}_{\rm def}Q\quad
 \forall x\in P:
           \|r_w(x)-e_w(h(x))\|<\delta.
 \tag{SF.3}
\]
Equivalently, the graph of $h$ must be a bisection of the coordinate
projections of the allowed relation
\[
 R_{w,\delta}=
   \{(x,y)\in P\times Q:
                    \|r_w(x)-e_w(y)\|<\delta\}.
 \tag{SF.4}
\]
A triangulation or a marked homeomorphism of this allowed set need not
preserve both coordinate projections. The argument for (SF.2) therefore
does not answer (SF.3).

A small explicit example isolates this loss of projection information.
Put $I=[0,1]$ and
\[
 q(x)=
 \begin{cases}
 2x,&0\le x\le1/3,\\
 1-x,&1/3\le x\le2/3,\\
 2x-1,&2/3\le x\le1.
 \end{cases}
 \tag{SF.5}
\]
The two formulas agree at both break points, $q(0)=0$, $q(1)=1$, and
$0<q(x)<1$ for $0<x<1$. The graphs of $\mathrm{id}_I$ and $q$ are
compact intervals with the same two marked endpoint positions. In fact
there is a semialgebraic homeomorphism of the ambient square, fixing its
entire boundary and preserving the first coordinate, which carries the
first graph onto the second. For $0<x<1$, set
\[
 H(x,y)=(x,h_x(y)),\qquad
 h_x(y)=
 \begin{cases}
 yq(x)/x,&0\le y\le x,\\
 q(x)+(y-x)(1-q(x))/(1-x),&x\le y\le1,
 \end{cases}
 \tag{SF.6}
\]
and define $H(0,y)=(0,y)$, $H(1,y)=(1,y)$. Each vertical map is strictly
increasing from zero to one. Its two expressions agree at $y=x$, with
value $q(x)$. Continuity on the sides $x=0,1$ follows directly: near
$x=0$, the lower part has both $x,y\to0$, while on the upper part
$q(x)\to0$ and $(1-q(x))/(1-x)\to1$; the argument at $x=1$ is the
same with the roles reversed. An explicit inverse, for $0<x<1$, is
\[
 h_x^{-1}(v)=
 \begin{cases}
 xv/q(x),&0\le v\le q(x),\\
 x+(1-x)(v-q(x))/(1-q(x)),&q(x)\le v\le1,
 \end{cases}
 \tag{SF.6a}
\]
with the identity on the sides $x=0,1$. These rational formulas are
continuous on their interior domains and agree at the dividing value.
For an explicit check at the sides, on the left-hand branch of $q$ the
forward displacement $h_x(y)-y$ equals $y$ for $y\le x$ and
$x(1-y)/(1-x)$ for $y\ge x$. On the right-hand branch it equals
$-(1-x)y/x$ for $y\le x$ and $y-1$ for $y\ge x$. Hence the supremum
displacements tend to zero as $x\to0$ and $x\to1$. The inverse has the
same supremum displacement, by substituting $v=h_x(y)$, so it is jointly
continuous at the sides as well. This proves that $H$ is a homeomorphism
over every real closed field, using no ordinary compactness argument.
It fixes all four sides.

Nevertheless $\Gamma(q)$ has no bisection covering both copies of $I$:
any section of its first projection is necessarily $x\mapsto(x,q(x))$,
and $q(1/4)=q(1/2)=q(3/4)=1/2$. The diagonal does have such a bisection.
The exact symbolic certificate is
\path{agents/maximal_locus/code/check_scaffold_projection_loss.py};
its executed output is
\path{agents/maximal_locus/logs/check_scaffold_projection_loss.json}.
It checks the breakpoint and endpoint identities, affine ranges, three
preimages, the forward and inverse vertical gluing and compositions,
positivity of their slopes, and the four side-displacement identities.
The continuity argument above supplies the topological part.

This example is a diagnostic for forgetting a projection. It is not a
counterexample to definability of (SF.3), since it is not being asserted
that every allowed relation of the special form (SF.4) has this graph
form. Nor is it a counterexample to A5.

\textbf{What a scaffold theorem would still have to prove.}
For A5.019 the unresolved witness is stronger than (SF.3): it is a
homeomorphism $A_t:L\times\Sigma\to L\times\Sigma$ with
\[
 \Theta_t A_t=\operatorname{pr}_{\Sigma},\qquad
 \sup_{(z,u)}\|\operatorname{pr}_L A_t(z,u)-z\|\longrightarrow0.
 \tag{SF.7}
\]
The first equation fixes a continuum of fibres. Replacing it by
compatibility with finitely many target strata or marked sets does not
retain that equation. The remaining finite-scaffold assertion would
need to extract, from the existence of an arbitrary witness of (SF.7),
one witness in a bounded collection of formula formats determined by
the fixed input family. Its bound must respect the exact angular
projection and the required error scale.

The already inspected Nayyeri--Xu proof bounds crossings for linear
genus-zero surfaces by convex shortcutting. A fixed nonlinear definable
triangulation does not preserve the cost inequalities used by that
shortcut. Refining a PL approximation as the tolerance tends to zero
does not supply a fixed formula bound. The self-contained folded-disk
extension proves the extension step once its affine boundary scaffold
exists; it does not construct a scaffold for curved angular data.
No extension of either result is assumed here.

Likewise A138's finite-format contractions and A160's relative
straightening theorem act on homeomorphisms already supplied in a
definable family. They can organize a bounded family of anchors once
it exists. Neither currently proves the witness extraction required in
(SF.7). This, rather than a continuity choice between already available
formulas, is the exact unresolved finite-scaffold step at this checkpoint.
